% Pursuing Stacks (cotes 134-2 to 134-8): archive cote 134-2, site page 43 (PDF page 44).
% appendix (letters to Breen). Draft compared: ps1.tex:1210-1212.
% AI-transcribed by Claude Opus 5.5 (claude-opus-5-5, adaptive thinking, effort high), Message Batches API, batch ps1 (msgbatch_01Rv5hCaNTtUt4JHx1yqoHPL), 2026-11-12; not yet checked by eye.
\scanpage{134-2/43}
\mspage{28}

O-simplexes, les 1-objets sont les chemins \xout{entre} (ou homotopies) entre
O-simplexes, les 2-objets sont les homotopies entre chemins (à extrémités fi-%
xées) etc. Pour les n-objets, cependant, on ne prend pas les homotopies
entre homotopies de fourbis, mais classes d'équivalence \xout{\unsure{entre}}\hand{\unsure{d'}}homotopies
(modulo la relation d'homotopie) entre homotopies. La composition des
i-objets ($i\hand{\geq}1$) se définit de façon évidente, on notera qu'elle n'est pas
strictement associative, mais associative modulo homotopie. Donc la n-caté-%
gorie qu'on obtient n'est pas"stricte" - et on prévoit pas mal d'emmerdement
pour définir de façon raisonnable une n-catégorie pas stricte (dans la
\unsure{déscription}\hand{/}des compatibilités \xout{entre} pour les "données d'associativité\hand{"}). La
mise sur pied du yoga qui suit pourrait constituer un fil d'\hand{A}riane pour la
définition en forme des n-c atégories (pas strictes), les n-foncteurs entre
elles (pas non plus stricts, et pour cause), les n-équivalences etc, au
même t\unsure{\xout{u}}itre que le yoga initial "une n-catégorie est une catégorie ou les
Hom et leurs accouplements\inserted{de composition} sont des(n-1)-catégories et des accouplements
entre telles". Cette n-catégorie\inserted{$C_n(K.)$} dépend fonctoriellement de $K.$ , tout
morphisme simplicial $K. \hand{$\to$} K'.$ définit un n-foncteur $C_n(K.)\hand{$\to$}C_n(K'.)$; en
fait, cela doit en dépendre même n-fonctoriellement, vu qu'on voit (en s'in%
spirant de ce qui précède et l'appliquant à des ensembles semi-simpliciaux
de la forme $Hom.(K.,K'.)$) que les ensembles semi-simpliciaux forment eux-%
mêmes \xout{une}\inserted{les O-objets d'une} n-catégorie, quel que soit $\unsure{N|n}$ ...

En fait, $C_n(K.)$ est un\xout{e} n-groupoïde, i.e. une n-catégorie où toute
i-fléche ($1\leq i\leq n$) (= i-objet) est une "équivalence" i.e. admet un quasi-%
inverse (donc un inverse si la n-catégorie est"réduite"). Si C est une
telle n-c atégorie i.e. un n-groupoïde, et x un O-objet de C, il s'impose
de désigner par $\hand{$\pi_i$}(C,x)$ ($O\leq i\leq n$) successivement: l'ensemble des classes
de O-objets à équivalence près (c'est un vulgaire ensemble), l'ensemble des
classes à équivalence près de 1-objets\hand{(}ou 1-flèches\hand{)} $x \to x$ (c'est un groupe,
pas nécessairement commutatif), l'ensemble des\inserted{classes mod équivalence des} 2-flèches $1_x \to 1_x$ , où $1_x$
est la 1-flèche identique de x (c'est un groupe commutatif $\hand{$\pi$}_2(C,x)$, ainsi
que les groupes qui vont suivre), l'ensemble des classes mod équivalence
de 3-flèches $1_{1_x} \longrightarrow 1_{1_x}$ , etc. Ces groupes forment, comme de juste, des
"systèmes loc aux" sur l'ensemble des O-objets de C, et modulo le grain de
sel habituel,\inserted{les $\pi_i(C,x)$} ne dépendent que de la"composante connexe" du O-objet x i.e.
de sa classe modulo équivalence de O-objets. Ceci dit, si C est de la form\unsure{e}
$C_n(K.)$, il résulte pratiquement des définitions que \xout{les} l'on a des isomor%
phismes canoniques $\hand{$\pi$}_i(K.,x)\hand{$\simeq$}\hand{$\pi$}_i(C_n(K.))$ pour $0\leq i\leq n$, qui pour x variable
peuvent s'interpréter comme des \over{\illegible{1}}{i}somorphismes de systèmes locaux. Il s'en-%
suit que pour une application semi-simpliciale \hand{$f$}:$K. \hand{$\to$} K.'$, le n-foncteur
correspondant $C_n(K.) \hand{$\to$} C_n(K.')$ est une n-équivalence sss f induit un iso-%

\margin{43}
